\section{Conclusion}
\label{sec:conclusion}

\paragraph{The answer to the guiding question.}
The question of Section~\ref{subsec:question} was why the inviscid Burgers'
equation needs an entropy condition and why that condition is the right one.
The chain of results is now complete. Smooth decreasing data develop an
infinite slope at an explicit time and place,
$T^*=-1/\min u_0'$ (Section~\ref{sec:shock-formation}); the weak formulation
continues the solution past this point through jumps obeying the
Rankine--Hugoniot condition (Section~\ref{sec:weak-solutions}), but it is
non-unique, already for the simplest Riemann data, where in fact infinitely
many weak solutions exist (Section~\ref{sec:nonuniqueness}). The Oleinik
entropy condition, of which the Lax condition $u_\ell>u_r$ is the local form
at a jump, restores uniqueness and gives the complete solution of the Riemann
problem (Sections~\ref{sec:entropy-condition}--\ref{sec:riemann-problem}).
That it is the \emph{right} condition is the content of the vanishing
viscosity limit: via the Cole--Hopf transformation the viscous solutions
converge to the entropy solution, whose jumps can only be downward because
the Lax--Oleinik minimiser is monotone, and directly because a viscous
travelling wave connecting $u_\ell$ to $u_r$ exists exactly when
$u_\ell>u_r$ (Section~\ref{sec:extensions}, Proposition~\ref{prop:viscous-profile}). The
theory itself is standard \cite[\S3.4]{evans2010}; the novelty of this
paper is in the explicit verifications and the computations below.

\paragraph{What the computations added.}
\begin{itemize}
  \item \emph{An explicit example} (Section~\ref{subsec:arctan-complete}): for
    $u_0=-\arctan x$ the entropy solution after breakdown has a single
    stationary shock of strength $\arctan y^*(t)$ ($y^*$ the positive root of
    $\arctan y=y/t$), growing like $\sqrt{3(t-1)}$ and saturating at $\pi/2$.
  \item \emph{A confirmed prediction}: a viscous shock of any strength has
    $L^1$ distance exactly $4\ln2\,\nu\approx2.77\,\nu$ from the sharp shock. The
    Cole--Hopf shock layer contributes $2.78\,\nu$ (within $0.5\%$), tending to
    $4\ln2\,\nu$ at every strength tested ($0.52$--$1.50$) once the window is
    small enough; the total error is first order, $\approx3.11\,\nu$, the
    remaining $\approx0.33\,\nu$ coming from the smooth part (measured, not
    derived).
  \item \emph{Godunov as a vanishing viscosity} (Section~\ref{subsec:godunov}):
    shock error $0.50\,\Delta x$ on aligned grids ($0.6$--$0.77\,\Delta x$
    otherwise); rarefaction order about $0.8$ and still rising. The modified
    equation predicts the shock error within $3\%$ at CFL $0.45$ (about $10\%$ over
    CFL $0.1$--$0.95$), i.e.\ $\nu_{\rm eff}\approx0.18\,\Delta x\to0$: a heuristic
    consistency check that a monotone scheme selects the entropy solution
    because its numerical diffusion acts as a vanishing viscosity; the
    rigorous result is \cite{crandall1980}.
  \item \emph{Corrected N-wave constants} (Section~\ref{subsec:n-wave}): $Q=-\min V$,
    $P=\int u_0-\min V$ with $V$ the primitive of $u_0$, not the positive and
    negative masses, since positive and negative parts cancel across shocks and
    only $P-Q=\int u_0$ is conserved. For $+1$ on $(0,1)$, $-1$ on $(1,2)$ the
    correct limit is $0$ ($\|u\|_{L^1}=1/t$ exactly), whereas the masses predict a
    wave of norm $2$; numerically only the corrected constants give decay.
\end{itemize}

\paragraph{Limits of the exposition.}
The proofs of the uniqueness theorem (Theorem~\ref{thm:oleinik-uniqueness}),
of the vanishing-viscosity limit (Theorem~\ref{thm:vanishing-viscosity}), of
the convergence of the Godunov scheme (Theorem~\ref{thm:godunov-convergence})
and of the N-wave decay (Theorem~\ref{thm:n-wave-decay}) are cited, not
reproduced. The numerical statements are observations for the specific data
sets and resolutions stated. Burgers' equation is a model of nonlinear
steepening and its regularisation, not a stepping stone to the far harder
open regularity questions for the three-dimensional Euler and Navier--Stokes
equations, whose mechanisms differ.
